\section{Introduction}\label{introduction}

What is non-Euclidean geometry? I believe most people in the modern
world have a vague understanding of what non-Euclidean means. It is
often connotated with a reality that isn't like ours and it is often
used as an umbrella term to describe all sorts of related concepts. In
this handout we try to give a mathematical introduction to the concepts
of non-Euclidean geometries. As a visual aide throughout this journey we
make use of tessellations.

This document serves as a self-study guide that picks up on topics from
discrete mathematics and linear algebra. The reader should be familiar
and comfortable handling graph theory, complex numbers and linear
transformations. It is worth looking up the following concepts as these
will serve as our basis.

\begin{itemize}
\tightlist
\item
  Euler Characteristic \(\chi = V - E + F\)
\item
  Euler Formula \(e^{i \phi} = \cos(\phi) + i \sin(\phi)\)
\item
  Polar coordinates \((\theta, r)\)
\item
  Spherical coordinates \((\theta, \phi, r)\)
\item
  Orthogonal Projection
\item
  Stereographic Projection
\end{itemize}

We highly encourage you to read further into the matter if you are
interested. This handout merely scrapes the top of this vast field of
mathematics.
\todo{mention notable mathematicians in this field? see (see tasks/260922-112332.md)}

\subsection{Historical Background}\label{historical-background}

In order for us to get a grasp on non-Euclidean geometries, we should
first understand what makes a geometry Euclidean in the first place. We
should all be relatively familiar with Euclidean geometry, it is what
most of us understand as \textbf{the} theory of geometry and what we are
taught out of high school. Euclid of Alexandria, our namesake,
formulated the ancient Greek's understanding of geometry, which later
formed the classical understanding of geometry. In his work Euclid gives
us the definition of geometric elements, which coincide with our modern
understanding. Furthermore Euclid states five Axioms.

\begin{enumerate}
\def\labelenumi{\arabic{enumi}.}
\tightlist
\item
  To draw a straight line {[}segment{]} from any point to any point.
\item
  To produce a finite straight line {[}segment{]} continuously in a
  straight line.
\item
  To describe a circle with any center and distance.
\item
  That all right angles are equal to one another.
\item
  That, if a straight line falling on two straight lines make the
  interior angles on the same side less than two right angles, the two
  straight lines, if produced indefinitely, meet on that side on which
  are the angles less than two right angles.{[}pp.~Euclid's elements{]}
\end{enumerate}

The first four axioms are very easily interpretable. Two points can be
connected by a straight line segment, line segments can be infinitely
extended into lines, there exists a circle with any center and radius
and all the right angles must be equal to one another.

The fifth postulate, also coined the parallel postulate, is where things
get interesting. At first glance it seems out of place. An simpler but
equivalent formulation this is given by Playfair's axiom. It states that
on a plane, given a line and a point not on it, only one parallel line
to the given line can be drawn through the point. Throughout the ages
mathematicians have speculated that the fifth axiom is superfluous and
could be derived from the first four. Many tried to find a proof without
success.

\missingfigure{illustration parallel postulate}

A geometry becomes non-Euclidean if we modify the fifth axiom
(\textbf{cite?}). There are two non-Euclidean geometries that arise from
this change. Elliptic Geometry, where there are exist zero parallels to
a given one and hyperbolic Geometry where there exist infinitely many.

\todo{When omiting the fifth axiom completely we talk about absolute geometry Faber 1983, pg. 131, }

\todo{State Hilbert's axioms, or mention it as the modern axiomatization, see handout from uni-bielefeld}

\subsection{Topology}\label{topology}

\subsection{Gaussian-Curvature}\label{gaussian-curvature}

When discussing non-Euclidean Geometry one will sooner or later fall
upon the term gaussian-curvature. Gaussian curvature is a measure of how
much a 2-dimensional surface curves at a specific point in space.

\todo{try to give analog from 1d curvature}

We call a curvature constant if it is the same at each point.

Gaussian-Curvature is an intrinsic property of a surface. Meaning that a
being living entirely on the surface could determine it's measure.

The sphere gives us a surface with constant positive curvature. We call
this surfaces embedded in \(\mathbb{R}^3\)

There exist surfaces with negative curvature in \(\mathbb{R}^3\), such
as the hyperboloid and paraboloid.

\todo{show parabolid and hyperbolid}

However Hilbert proofed that there are no embeddings of a surface with
constant negative curvature in 3-d space. This is the main reason why it
is so hard for us to imagine hyperbolic space.

\subsection{Tessellations}\label{tessellations}

A \emph{tessellation} (or Tiling) is the process of covering a surface
with geometrical shapes so that there are no gaps or overlaps. The
shapes used in a tiling are called \emph{tiles}. Tilings are classified
by restricting the choice and alignment of tiles.

\subsubsection{Polygonal Tessellations}\label{polygonal-tessellations}

A \emph{polygonal} tessellation is a tessellation where we restrict the
tiles to be polygons. If we further restrict the alignment of tiles,
such that adjacent tiles share a common edge and meet at their vertices,
we talk about an \emph{edge-to-edge} tiling.

Polygonal edge-to-edge tilings can be classified further using the
following properties.

A tiling is said to be \emph{regular}, if all it's tiles are regular
polygons.

A tiling is said to be \emph{face}-transitive, if all it's tiles are
congruent.

A tiling is said to be \emph{edge}-transitive, if

\subsubsection{Triangle Tiling}\label{triangle-tiling}
